% ---------------------------------------------------------------------
\subsection{The simulated room: surfaces and early reflections (\fn{RoomSim.h})}
\label{sec:math-saal}
The \ui{Saal} (\S\ref{sec:room}) is computed from geometry. Coordinates are metres with $x$ to the right (0 at the left
wall), $y$ to the front (0 at the back wall), $z$ up (0 at the floor); the listener faces $+y$ at the ear height of
\SI{1.6}{m} (or half the height in a room lower than \SI{3.2}{m}).

\paragraph{Absorption in two bands (\fn{SurfAbs}).} Every surface has an energy absorption coefficient $\alpha$ in a
low/mid band and a high band. A surface that follows the sliders takes $\alpha_\mathrm{lo}$ from \ui{Absorption Wände}
or \ui{Boden/Decke}; a surface with a material of its own takes it from the table (stone 0.02, plaster 0.05, wood 0.12,
glass 0.04, curtain 0.5, mirror 0 as a wall). The high band adds a share of what is left,
\begin{equation}
  \alpha_\mathrm{hi} = \min\bigl(0.99,\ \alpha_\mathrm{lo} + 0.6\,\eta\,(1-\alpha_\mathrm{lo})\bigr),
\end{equation}
with the \ui{Höhendämpfung} $\eta$ (or the material's own value): highs are always absorbed at least as much as
lows, as by every real material. Air absorbs energy at the rate $m$ per metre, $m_\mathrm{lo} = 0.0002$ and
$m_\mathrm{hi} = 0.004 + 0.012A$ with the \ui{Luft} amount $A$.

\paragraph{Image sources (\fn{imageSources}).} In a rectangular room a reflection off a wall is the same as a straight
path from the source's mirror image behind that wall \citep{allen1979image}. Along one axis of length $D$ with the
source at $s$, the image after $m$ reflections lies at
\begin{equation}
  x_m = mD + \begin{cases} s, & m \text{ even}\\ D - s, & m\text{ odd}\end{cases},\qquad m\in\mathbb Z,
\end{equation}
(\fn{imageCoord}); an even number of mirrorings restores the orientation, an odd number flips it. A three-dimensional
image is a triple $(m_x,m_y,m_z)$ with $|m_x|+|m_y|+|m_z| = k \le K$ reflections, $K\le3$: 6 images of order 1, 18 of
order 2 and 38 of order 3 (62 in all). For each axis the image's path hits the wall at 0 $\lceil |m|/2\rceil$ times
and the wall at $D$ $\lfloor |m|/2\rfloor$ times if $m<0$, the other way round if $m>0$ (\fn{axisHits}), so every hit
can be charged to the right surface.

\paragraph{The gain of a reflection (\fn{updateSlotTaps}).} An image at distance $r$ from the listener, whose path hits
surfaces $i_1,\dots,i_k$, is heard after the delay $(r - r_d)/c$ relative to the direct sound ($r_d$ the direct distance;
the dry pipe is already the direct sound, so all reflections are delayed relative to it) with the amplitude
\begin{equation}
  g = g_\mathrm{base}\, g_d(r)\, g_\mathrm{ER}\,\prod_{j=1}^{k}\sqrt{1-\alpha_{\mathrm{lo},i_j}}\;\cdot\;\bigl(\sqrt{1-\sigma}\bigr)^{k},
\end{equation}
where $\sqrt{1-\alpha}$ is the pressure reflection factor of a surface that absorbs the energy share $\alpha$, $g_d$ the
distance law of \S\ref{sec:math-motion}, and $\sigma$ the \ui{Streuung}: the share $\sigma$ of the energy is scattered at
every hit and goes to the diffuse tail instead, so the specular image keeps $1-\sigma$ of the energy per reflection. The
highs are treated as a shelf relative to the lows,
\begin{equation}
  \rho = \prod_j \sqrt{\frac{1-\alpha_{\mathrm{hi},i_j}}{1-\alpha_{\mathrm{lo},i_j}}}\;\cdot\;e^{-\frac12(m_\mathrm{hi}-m_\mathrm{lo})\,r},
\end{equation}
(the second factor the extra air absorption of the highs over the path, as an amplitude: half the energy rate). Each tap
is rendered as $y = \mathrm{LP}(x) + \rho\,(x - \mathrm{LP}(x))$ with a one-pole low-pass at \SI{2.5}{kHz}: the band
below passes with the gain $g$, the band above with $g\rho$. Its direction is encoded with \fn{encodeAmbiXDir} from the
vector to the image, $(\cos\varphi,\sin\varphi) = (\Delta y, -\Delta x)/\rho_{xy}$ (the listener faces $+y$ and
azimuth is positive to the left), at up to third order for first-order reflections and at first order for later ones,
which are many, quieter and more diffuse. Reflections later than the delay buffer are faded out linearly over its last
15\,\% (they belong to the tail anyway).

\paragraph{Moving sources.} The delay of a tap is read from a ring buffer with linear interpolation,
$x = b + f(a-b)$ between the two neighbouring samples, and ramped linearly across the block from the old delay $d_0$ to
the new one $d_1$, like the gains. A delay that changes smoothly is the Doppler shift of the reflection: every image
moves when the pipe moves and its pitch changes with its own radial speed. The integer and fractional parts of the
read position are kept apart (an integer clock and a float offset), because a single float clock would lose the
fraction after a few minutes of running and alias. A tap whose inputs did not change is not recomputed: the key of the
computation (source direction and distance, the room's dimensions and materials, the listener and the order) is
compared with the stored one.

\paragraph{Polygon rooms.} In a floor plan that is not a rectangle the images are found by mirroring over the wall
lines in the plan, $\mathbf p' = \mathbf p - 2(\mathbf n\cdot\mathbf p + c)\,\mathbf n$ for a wall with inward normal
$\mathbf n$ and offset $c$ (\fn{reflectOver}), for every sequence of up to three walls, combined with the floor and
ceiling images in $z$. A path is valid only if every reflection point lies on its wall segment and, in a non-convex
plan, if no other wall blocks one of its legs (\fn{planPathPoints}, \fn{planBlocked}). Since this search is expensive
it runs only when the source moved noticeably; in between, each tap keeps its path and only recomputes its image
position, and paths that appear or disappear fade in or out over \SI{25}{ms}, like a reflection sliding past the end of
a wall.

\paragraph{Around a corner (\fn{diffractedPath}, \fn{diffractionLoss}).} When a wall of the plan lies between pipe and
listener, the direct sound takes the shortest way around one or two reflex corners of the plan (moved \SI{5}{cm} into
the room) that both can see. It is heard from the direction of the first corner, from the distance of the whole way,
and loses level by Maekawa's rule for the insertion loss of a barrier \citep{maekawa1968}, from the Fresnel number
$N = 2\delta/\lambda$ of the detour $\delta$ (the way around minus the straight line), evaluated at \SI{250}{Hz}:
\begin{gather}
  L_\mathrm{loss} = \min\Bigl(26,\ 10\log_{10}(1+20N) + 4b\Bigr)\ \mathrm{dB},\qquad
  N = \frac{2\delta\cdot\SI{250}{Hz}}{c},\\
  f_c = 250\Bigl(4 + \frac{3}{29.2\,\delta}\Bigr)\,\mathrm{Hz},
\end{gather}
with the bend $b\in[0,1]$ ($0$--$180^\circ$) adding a little for a sharp turn. Above $f_c$, where the loss would be
\SI{6}{dB} higher, the low-pass of \S\ref{sec:math-voice} sets in, so the highs are lost first, as the detour counts in
wavelengths. At the edge of the shadow $\delta = 0$ and nothing is lost: stepping past a corner never makes the sound jump.

% ---------------------------------------------------------------------
\subsection{The simulated room: decay and tail (\fn{RoomSim.h}, \fn{RoomTail.cpp})}
\label{sec:math-tail}

\paragraph{Eyring's formula (\fn{eyringRt}).} For a quick estimate (the page shows it) the reverberation time per band is
\begin{equation}
  T_{60} = \frac{0.161\,V}{-S\ln(1-\bar\alpha) + 4mV},\qquad \bar\alpha = \frac{\sum_i\alpha_iS_i}{S},
\end{equation}
with the volume $V$, the total surface $S$ (in a polygon room the plan's area twice and the walls' lengths times the
height) and the area-weighted mean absorption \citep{eyring1930}. $-\ln(1-\bar\alpha)$ instead of Sabine's
$\bar\alpha$ accounts for the fact that energy is lost in discrete reflections, not continuously, which matters in
absorbent rooms; $4mV$ is the air. The constant $0.161 = 24\ln 10/c$ (in s/m) comes from $c = \SI{343}{m/s}$. The
mean free path, used for the onset of the tail and the FDN, is $\bar\ell = 4V/S$.

\paragraph{Ray tracing (\fn{traceRoom}).} Eyring assumes a diffuse field and evenly spread absorption, which a long nave
with a dead floor and hard walls violates. So the decay actually used is measured by tracing 3000--4000 rays from
where an organ would stand ($W/2$, $3L/4$, $\min(H/2, \SI{4}{m})$) in uniformly random directions ($z$ uniform in
$[-1,1]$, azimuth uniform, which is uniform on the sphere by Archimedes' theorem). Each ray carries an energy per band,
starting at 1. Between two hits it travels $\ell$ metres in $\ell/c$ seconds and deposits its energy times the time
spent into bins of \SI{5}{ms}; at the end of the segment it loses $e^{-m\ell}$ to the air and a factor $1-\alpha$ of the
surface it hits. It then continues specularly ($\mathbf d \leftarrow \mathbf d - 2(\mathbf d\cdot\mathbf n)\mathbf n$),
or with the probability $\sigma$ diffusely, in a cosine-weighted (Lambert) direction around the normal,
\begin{equation}
  \mathbf d = \sqrt{u_1}\cos(2\pi u_2)\,\mathbf t_1 + \sqrt{u_1}\sin(2\pi u_2)\,\mathbf t_2 + \sqrt{1-u_1}\,\mathbf n,
\end{equation}
which draws directions with a density proportional to $\cos\theta$ (Malley's method: uniform on the disc, lifted to the
hemisphere). A ray ends after 4000 hits, \SI{30}{s}, or when both energies are below $10^{-7}$. The same seed is used
every time, so modulating the room does not make the decay jump randomly.

\paragraph{From the energy to $T_{30}$ (\fn{fitDecay}).} The summed energy per bin $e_b$ is integrated backwards,
$S_b = \sum_{b'\ge b} e_{b'}$ (Schroeder's backward integration, which turns the noisy ensemble decay into a smooth,
monotone curve), converted to $\mathrm{dB} = 10\log_{10}(S_b/S_0)$, and a straight line is fitted by least squares
to the part between $-5$ and \SI{-35}{dB}. With the slope $k$ in dB per second, $T = -60/k$: the $T_{30}$ of room
acoustics, extrapolated to 60 dB.

\paragraph{The synthetic tail (\fn{TailIR::build}).} The late reverberation is a response synthesised from the two
decay times $T_\mathrm{lo}$, $T_\mathrm{hi}$ ($T_\mathrm{hi}\le T_\mathrm{lo}$ enforced). For every output channel
$c$ (up to second order: a diffuse field needs no more) a noise $\nu_c$ is made from four uniform numbers,
$\nu = 0.866\sum_{i=1}^4 b_i$, which by the central limit theorem is close to Gaussian with unit variance
($\operatorname{Var}(b) = 1/3$, so $4\cdot\tfrac13\cdot0.866^2 \approx 1$). It is split at \SI{1500}{Hz} by a biquad
low-pass into $\nu_\mathrm{lo}$ and $\nu_\mathrm{hi} = \nu-\nu_\mathrm{lo}$ (the two always sum to the original) and each
band decays with its own time:
\begin{equation}
  h_c(t) = w_l\,\Bigl(\nu_{\mathrm{lo}}(t)\,e^{-6.9078\,t/T_\mathrm{lo}} + \nu_\mathrm{hi}(t)\,e^{-6.9078\,t/T_\mathrm{hi}}\Bigr),
  \qquad t \ge t_0,
\end{equation}
where $6.9078 = \ln 10^3$ makes the amplitude fall by $\SI{60}{dB}$ in $T$. The tail starts at $t_0 = \bar\ell/(2c)$,
after the first reflections, and fades in by a raised cosine over one mean free time $\bar\ell/c$. The channel weight
$w_l = 1/\sqrt{2l+1}$ gives a diffuse field its SN3D energy ratios, as in \S\ref{sec:math-ir}; the seed is fixed per
channel, so the channels are uncorrelated with each other and the same from one rebuild to the next (a modulated room
changes its envelopes, not its noise, and a held chord's reverb does not jump). The response is made long enough to
decay by \SI{75}{dB}; if that would exceed the maximum length, the decay is bent smoothly: from 35\,\% of the length on, an
extra loss growing with the square of the time, $\Delta(u) = (75 - 60\,t_\mathrm{len}/T)\,u^2$ dB with $u$ the
fraction of the remaining span, brings it to $-75$ dB exactly at the end without a corner. The response is normalised to
unit energy in $W$ and run through the three-stage partitioned convolution of \S\ref{sec:math-partconv}
(blocks 256, 2048, 8192).

\paragraph{The feedback delay network (\fn{renderFdn}).} The earlier tail, still selectable (\ui{FDN}), is a recursive
reverberator \citep{jot1991fdn} with 16 delay lines of lengths
\begin{equation}
  d_j = \frac{\bar\ell}{c}\,f_s\,\rho_j,\qquad \rho_j\in\{0.53, 0.59, 0.67, \dots, 1.41, 1.49\},
\end{equation}
spread around the mean free time with mutually prime-like ratios, so their echoes do not line up. Each line's output
is split at \SI{1500}{Hz} by a one-pole low-pass and each band is attenuated by the gain that makes it decay with its
band's ray-traced time,
\begin{equation}
  g_{j} = 10^{-3 d_j / (f_s T)},
\end{equation}
(60 dB per $T$, distributed over the line's length: a signal that has circulated for $T$ seconds has passed through
gains whose product is $10^{-3}$). The sixteen outputs are mixed by a fixed random orthogonal matrix $\mathbf A$
(16 random rows made orthonormal by Gram--Schmidt): an orthogonal matrix preserves energy, so the loop's decay is set
by the gains alone, and a dense random one mixes every line into every other at once (a Hadamard matrix left isolated
modes around 1.5--\SI{2}{kHz}). The input, delayed by half a mean free time, enters every line with $\pm0.25$. Every
delay is modulated by $\pm8$ samples (at \SI{48}{kHz}) with its own rate in $[0.3,1.2]$\,Hz, computed with a rotating
phasor $(c,s)\leftarrow(c\cos\omega - s\sin\omega,\ s\cos\omega + c\sin\omega)$ that is renormalised once per block,
which smears the modes so no frequency rings. Line $j$ is heard from a direction of a 16-point Fibonacci sphere
(azimuth $137.508^\circ j$, the golden angle, $\sin\vartheta = 1 - 2(j+\frac12)/16$), so the tail surrounds the
listener. The output level is $\sqrt{1-\bar g^2}$ times a constant, which keeps the reverb's energy roughly independent
of the decay time ($\sum_n \bar g^{2n} = 1/(1-\bar g^2)$). When the room changes, the delays, the pre-delay and the gain
glide with a time constant of \SI{80}{ms}.
